%% ma: L10-B15-0048
%% lop: 10
%% chuong: 6
%% bai: 15
%% dang: 10.15.5
%% loai: TL
%% muc: TH
%% dap_an: "-"
%% nguon: "Ảnh giáo viên gửi 21/11/2026: Chuyên đề VI – Toán 10 – Chương VI (Hàm số – đồ thị và ứng dụng), tr. 199–203, Dạng 5 Câu 1"
%% kien_thuc: "Lớp 10 Bài 15 (hàm số đồng biến, nghịch biến)"
%% trang_thai: chinh-thuc
%% ghi_chu: "Dạng do giáo viên cung cấp (21/11/2026)."
\begin{ex}
Có bao nhiêu giá trị nguyên của tham số $m$ thuộc đoạn $[-3;3]$ để hàm số $f(x)=(m+1)x+m-2$ đồng biến trên $\mathbb R$?
\loigiai{$\dfrac{f(x_1)-f(x_2)}{x_1-x_2}=m+1$. Hàm số đồng biến $\Leftrightarrow m+1>0\Leftrightarrow m>-1$. Với $m\in[-3;3]$ nguyên: $m\in\{0;1;2;3\}$. Có $4$ giá trị.}
\end{ex}
